% !TeX root = ../main.tex

\section{การทดลองที่ 5: 2-State KF --- ออกแบบ, Implement, และทดสอบ Free Vibration}
\textit{ครอบคลุม Criteria: Part 3 ทั้งหมด (รวม 9 คะแนน) --- State-Space Model,
Discretization, เหตุผลการเลือก Sampling Rate, เหตุผลการใช้ Equilibrium เป็น State,
เขียนโปรแกรม 2-State KF, หา Matrix $Q$ และ $R$, อธิบายความหมายทางกายภาพของ $Q$
และ $R$, ทดสอบระยะต่างๆ, เปรียบเทียบกับ Model, วิเคราะห์ Velocity Estimate, สรุปผล}

%% ─────────────────────────────────────────────────────────────────────────────
\subsection{บทนำ}

การทดลองที่ 1--4 ได้พารามิเตอร์ครบสำหรับการสร้างแบบจำลองระบบ Mass-Spring-Damper
ได้แก่ $K = 29.11$~N/m จากการทดลองที่ 1, มวล $m$ จากการชั่ง,
และ Damping Coefficient $C$ จาก Parameter Estimator ในการทดลองที่ 4
การทดลองที่ 5 จึงเป็นการนำพารามิเตอร์เหล่านี้ทั้งหมดมาบูรณาการ
เพื่อสร้าง \textbf{2-State Discrete Kalman Filter} ที่สมบูรณ์

เป้าหมายหลักของ 2-State KF คือการประมาณค่า State สองตัวพร้อมกัน ได้แก่
\textbf{Position} $x_1$ ที่วัดได้โดยตรงจาก HC-SR04 แต่มี Noise
และ \textbf{Velocity} $x_2$ ที่ \textit{ไม่สามารถวัดได้โดยตรง} ด้วย Sensor ที่มี
ความสามารถนี้คือสิ่งที่ทำให้ Kalman Filter มีคุณค่าในระบบควบคุมจริง
เพราะระบบสามารถ``รู้''ความเร็วของมวลได้โดยไม่ต้องติดตั้ง Velocimeter เพิ่ม

ในการทดลองนี้จะพิสูจน์ว่า 2-State KF สามารถ (1) ประมาณ Position ได้เรียบกว่า
Raw Sensor Data, (2) ประมาณ Velocity ได้อย่างสมเหตุสมผลทางกายภาพ,
และ (3) ให้ผลที่สอดคล้องกับ Mathematical Model ที่ได้จากการทดลองที่ 4

%% ─────────────────────────────────────────────────────────────────────────────
\subsection{วัตถุประสงค์}

\begin{itemize}
    \item เพื่อสร้าง 2-State (Position + Velocity) Discrete Kalman Filter
          สำหรับระบบ Mass-Spring-Damper บน STM32 Nucleo G474RE
    \item เพื่อ Discretize State-Space Model ($A_c$, $B_c$, $H$)
          และเลือก Sampling Rate ที่เหมาะสมพร้อมเหตุผลทางทฤษฎี
    \item เพื่อกำหนดค่า Matrix $Q$ (2$\times$2) และ Scalar $R$
          พร้อมอธิบายความหมายทางกายภาพของแต่ละ Element
    \item เพื่อทดสอบ KF ด้วย Free Vibration ที่ Initial Displacement
          ต่างกัน 3 ระดับ และเปรียบเทียบกับ Mathematical Model
    \item เพื่อวิเคราะห์ว่า Velocity Estimate ที่ KF คำนวณได้
          มีความสมเหตุสมผลทางกายภาพหรือไม่
\end{itemize}

%% ─────────────────────────────────────────────────────────────────────────────
\subsection{สมมติฐาน}

\begin{itemize}
    \item 2-State KF จะให้ค่าประมาณ Position ที่เรียบกว่า Raw Data
          โดยยังสอดคล้องกับ Mathematical Model ที่ Identify ได้
    \item Velocity Estimate จาก KF (ที่ไม่ได้วัดโดยตรง) จะมีความสอดคล้อง
          ทางกายภาพกับ Derivative ของ Displacement โดยเฉพาะอย่างยิ่ง
          Velocity ควรเป็นศูนย์ที่จุด Maximum Displacement
          และสูงสุดที่จุด Equilibrium Crossing
    \item $Q$ ที่ใหญ่ขึ้นจะทำให้ KF ติดตาม Measurement มากขึ้น
          ส่วน $R$ ที่ใหญ่ขึ้นจะทำให้ KF พึ่ง Model Prediction มากขึ้น
    \item ที่ Displacement ขนาดใหญ่ ($x_0 = 6$~cm) Linear Model
          อาจเบี่ยงเบนจาก Response จริงเนื่องจากความไม่เชิงเส้นของสปริง
          และแรงเสียดทาน
\end{itemize}

%% ─────────────────────────────────────────────────────────────────────────────
\subsection{ตัวแปร}

\begin{itemize}
    \item \textbf{ตัวแปรอิสระ:} Initial Release Displacement $x_0$
          ที่ 2~cm, 4~cm, และ 6~cm; ค่า Matrix $Q$ ที่ปรับเปลี่ยนใน
          Sensitivity Study; ค่า Scalar $R$ ที่ปรับเปลี่ยน
    \item \textbf{ตัวแปรตาม:} ค่าประมาณ Position $\hat{x}_1$ (m),
          ค่าประมาณ Velocity $\hat{x}_2$ (m/s),
          Estimation Error ระหว่าง KF Output กับ Raw Measurement,
          ความสอดคล้องระหว่าง KF Output กับ Mathematical Model
    \item \textbf{ตัวแปรควบคุม:} ค่า $m$, $K = 29.11$~N/m, $C$
          คงที่จากการทดลองที่ 1 และ 4, Sampling Time $T_s$ คงที่,
          Control Input $u = 0$ (Free Vibration)
\end{itemize}

%% ─────────────────────────────────────────────────────────────────────────────
\subsection{เอกสารและงานวิจัยที่เกี่ยวข้อง}

\subsubsection{ค่า Spring Constant ที่ใช้ในแต่ละการทดลอง}

เนื่องจากการทดลองที่~4 และ~5 ใช้ Setup ที่เพิ่ม Guide Rail
เข้ามาเพื่อบังคับให้มวลเคลื่อนที่ในแนวดิ่งอย่างเคร่งครัด
ทำให้ค่า $K$ ที่ Estimate ได้ต่างจากการทดลองที่~1
ตารางที่~\ref{tab:K_summary} สรุปค่า $K$ ที่ใช้ในแต่ละการทดลอง

\begin{table}[htbp]
\centering
\caption{ค่า Spring Constant $K$ ที่ใช้ในแต่ละการทดลอง}
\label{tab:K_summary}
\begin{tabular}{clcc}
\toprule
การทดลอง & วิธีได้มา & $K$ (N/m) & เหตุผล \\
\midrule
1 & Full OLS Regression & 29.11 & Setup ไม่มี Guide Rail \\
4 & MATLAB Parameter Estimator & 32.04 & Setup มี Guide Rail \\
5 & ใช้ค่าจากการทดลองที่~4 & 32.04 & ใช้ Setup เดียวกัน \\
\bottomrule
\end{tabular}
\end{table}

\subsubsection{State-Space Representation ของ Second-Order System}

สมการ EOM ของ MSD (สมการ~\eqref{eq:eom_msd} จากการทดลองที่ 4)
เป็น Second-Order ODE ซึ่งสามารถแปลงเป็น First-Order State-Space Form
โดยกำหนด State Vector $\mathbf{x} = [x_1,\; x_2]^T = [\text{position},\; \text{velocity}]^T$:

\begin{equation}
    \dot{x}_1 = x_2
    \label{eq:ss_x1dot}
\end{equation}
\begin{equation}
    \dot{x}_2 = -\frac{K}{m} x_1 - \frac{C}{m} x_2
    \label{eq:ss_x2dot}
\end{equation}

เหตุผลที่เลือก Position และ Velocity เป็น State Vector คือทั้งสองตัวแปร
ร่วมกันอธิบายพลศาสตร์ของระบบได้ครบถ้วน (ทราบทั้งสองแล้วสามารถทำนาย
พฤติกรรมในอนาคตได้ทั้งหมด) ซึ่งเป็นข้อกำหนดพื้นฐานของ State Variable

\subsubsection{Continuous-Domain State-Space Matrices}

จากสมการ~\eqref{eq:ss_x1dot}--\eqref{eq:ss_x2dot} เขียนในรูป Matrix ได้:

\begin{equation}
    \dot{\mathbf{x}}(t) = A_c \mathbf{x}(t) + B_c u(t), \qquad
    z(t) = H\mathbf{x}(t) + v(t)
    \label{eq:ss_compact}
\end{equation}

โดย

\begin{equation}
    A_c = \begin{bmatrix} 0 & 1 \\ -K/m & -C/m \end{bmatrix}, \quad
    B_c = \begin{bmatrix} 0 \\ 1/m \end{bmatrix}, \quad
    H   = \begin{bmatrix} 1 & 0 \end{bmatrix}
    \label{eq:cont_matrices}
\end{equation}

\subsubsection{เหตุผลที่ $H = [1,\; 0]$}

HC-SR04 วัดได้เฉพาะ Position ($x_1$) ไม่สามารถวัด Velocity ($x_2$) ได้โดยตรง
ดังนั้น Observation Matrix จึงต้อง``เลือก''เฉพาะ Element แรกของ State Vector:

\begin{equation}
    z_k = H\mathbf{x}_k + v_k = \begin{bmatrix}1 & 0\end{bmatrix}
          \begin{bmatrix}x_1 \\ x_2\end{bmatrix} + v_k = x_1 + v_k
    \label{eq:obs_equation}
\end{equation}

KF จัดการกับ \textbf{Partial Observation} นี้โดยใช้ $A_d$ เพื่อ Propagate
ข้อมูล Velocity ผ่านรอบเวลา Prediction Step จะแพร่ข้อมูล Position
ไปยัง Velocity ผ่านองค์ประกอบ $(A_d)_{21}$ ทำให้แม้จะวัดได้แค่ Position
KF ก็สามารถ Estimate Velocity ได้ผ่านกระบวนการ Predict-Update Loop

\subsubsection{เหตุผลที่ $u = 0$ ใน Free Vibration}

จากการอธิบายใน Equilibrium Cancellation ของการทดลองที่ 4:
เมื่อใช้การกระจัดจาก Equilibrium เป็น State แรงโน้มถ่วง $mg$
ถูก Balance โดย Spring Pre-load $Kx_\text{eq} = mg$ พอดี
ทำให้ EOM กลายเป็น $m\ddot{x} + C\dot{x} + Kx = 0$
ซึ่งไม่มี External Input ดังนั้น $u = 0$ และ $B_d = \mathbf{0}$
ตลอดการทดลอง Free Vibration

\subsubsection{การเลือก Sampling Rate}

ตาม Nyquist-Shannon Sampling Theorem Sampling Frequency $f_s$
ต้องสูงกว่า 2 เท่าของ Natural Frequency สูงสุดในระบบ
แต่สำหรับ Control System ในทางปฏิบัติควรใช้ 10--20 เท่าเพื่อความเสถียร:

\begin{equation}
    f_s \geq 20\, f_n, \qquad \text{หรือ} \qquad
    T_s \leq \frac{1}{20\, f_n} = \frac{2\pi}{20\,\omega_n}
    \label{eq:sampling_criterion}
\end{equation}

โดยที่ $f_n = \omega_n / (2\pi)$ คือ Natural Frequency ในหน่วย Hz
เหตุผลที่ใช้ 20 เท่าแทน 2 เท่า เพราะต้องการให้สัญญาณที่ Discretize แล้ว
ยังคงจำลองพลศาสตร์ Continuous ได้อย่างแม่นยำ
และ Euler Forward Approximation จะมี Truncation Error น้อยลง

\subsubsection{Discretization: Euler Forward vs.\ ZOH}

\paragraph{Euler Forward (First-Order Approximation)}

\begin{equation}
    A_d = I + A_c T_s =
    \begin{bmatrix} 1 & T_s \\ -\dfrac{K}{m}T_s & 1 - \dfrac{C}{m}T_s \end{bmatrix}
    \label{eq:euler_Ad}
\end{equation}

\begin{equation}
    B_d = B_c T_s = \begin{bmatrix} 0 \\ T_s/m \end{bmatrix}
    \label{eq:euler_Bd}
\end{equation}

วิธีนี้ง่ายในการ Implement แต่มี Truncation Error $O(T_s^2)$
ที่ $T_s = 10$~ms ค่า $A_d$ จาก Euler Forward ต่างจาก ZOH ประมาณ
$2\%$ ซึ่งยอมรับได้สำหรับ Application นี้

\paragraph{Zero-Order Hold (ZOH)}

\begin{equation}
    A_d = e^{A_c T_s} = I + A_c T_s + \frac{(A_c T_s)^2}{2!} + \cdots
    \label{eq:zoh_Ad}
\end{equation}

วิธีนี้แม่นยำกว่าและรักษา Eigenvalue ของระบบ Continuous ไว้ได้ดีกว่า
แต่ต้องคำนวณ Matrix Exponential ซึ่งซับซ้อนกว่าสำหรับ Embedded System

ทั้งสองวิธีให้ผล $A_d$ ที่ Stable (Eigenvalue อยู่ภายใน Unit Circle)
สำหรับพารามิเตอร์ในการทดลองนี้

\subsubsection{ความหมายทางกายภาพของ Matrix $Q$ และ Scalar $R$}

\paragraph{Measurement Noise Covariance $R$ (Scalar)}

$R$ คือ Variance ของ Noise ของ HC-SR04 ที่คำนวณจากการทดลองที่ 2:
\begin{equation}
    R = \sigma^2_\text{sensor} = 0.0337 \; \text{cm}^2
    \label{eq:R_exp5_nominal}
\end{equation}

$R$ บอกว่า Filter ควรเชื่อถือ Measurement มากแค่ไหน
ค่าสูงหมายความว่า Sensor Noisy มาก Filter จะพึ่ง Model Prediction มากขึ้น

\paragraph{Process Noise Matrix $Q$ (2$\times$2)}

$Q$ เป็น Diagonal Matrix ที่แสดงความไม่แน่นอนในแต่ละ State:
\begin{equation}
    Q = \begin{bmatrix} q_1 & 0 \\ 0 & q_2 \end{bmatrix}
    \label{eq:Q_structure}
\end{equation}

ความหมายทางกายภาพของแต่ละ Element:
\begin{itemize}
    \item \textbf{$q_1 = Q_{11}$ (Position Process Noise):}
          ความไม่แน่นอนในการทำนาย Position ในแต่ละ Step
          ค่านี้มักมีขนาดเล็กเพราะ Position เปลี่ยนแปลงช้า
          (เป็น Integral ของ Velocity) และ Dynamic Model
          อธิบาย Position ได้แม่นยำในช่วง Displacement เล็ก
    \item \textbf{$q_2 = Q_{22}$ (Velocity Process Noise):}
          ความไม่แน่นอนในการทำนาย Velocity ในแต่ละ Step
          ค่านี้มักใหญ่กว่า $q_1$ เพราะ Velocity ได้รับผลกระทบ
          จากแรงที่ Model ไม่ได้รวม เช่น แรงเสียดทาน Non-linear
          และ Air Resistance ซึ่งทำให้ Model Prediction ของ Velocity
          มีความคลาดเคลื่อนสะสมมากกว่า
    \item \textbf{Off-diagonal elements = 0:}
          สมมติว่า Process Noise ของ Position และ Velocity
          ไม่มี Correlation กัน (Statistically Independent)
          ซึ่งเป็นการประมาณที่สมเหตุสมผลและลดความซับซ้อนในการ Tune
\end{itemize}

%% ─────────────────────────────────────────────────────────────────────────────
\subsection{ทฤษฎี: 2-State Discrete Kalman Filter สำหรับ MSD}

\subsubsection{Continuous State-Space Model}

กำหนด State Vector:
\begin{equation}
    \mathbf{x}(t) = \begin{bmatrix} x_1(t) \\ x_2(t) \end{bmatrix}
                  = \begin{bmatrix} \text{Displacement จาก Equilibrium (m)} \\
                                    \text{Velocity (m/s)} \end{bmatrix}
    \label{eq:state_vector}
\end{equation}

จาก EOM ของ MSD ในสมการ~\eqref{eq:eom_msd} แปลงเป็น First-Order Form:
\begin{equation}
    \dot{\mathbf{x}}(t) = A_c \mathbf{x}(t) + B_c u(t), \qquad
    z(t) = H \mathbf{x}(t) + v(t)
    \label{eq:ss_cont}
\end{equation}

\begin{equation}
    A_c = \begin{bmatrix} 0 & 1 \\ -K/m & -C/m \end{bmatrix}, \quad
    B_c = \begin{bmatrix} 0 \\ 1/m \end{bmatrix}, \quad
    H   = \begin{bmatrix} 1 & 0 \end{bmatrix}
    \label{eq:ss_matrices}
\end{equation}

เมื่อแทนค่าตัวเลขพารามิเตอร์จริงที่ได้จากการทดลอง:
\begin{equation}
    A_c = \begin{bmatrix} 0 & 1 \\ -K/m & -C/m \end{bmatrix}
    = \begin{bmatrix} 0 & 1 \\ [\text{ค่าจริง}] & [\text{ค่าจริง}] \end{bmatrix}
    \label{eq:Ac_numerical}
\end{equation}

\subsubsection{การเลือก Sampling Time}

คำนวณ Natural Frequency จากพารามิเตอร์ที่ได้:
\begin{equation}
    \omega_n = \sqrt{\frac{K}{m}}, \qquad f_n = \frac{\omega_n}{2\pi}
    \label{eq:omega_n}
\end{equation}

ใช้เกณฑ์ 20 เท่าของ $f_n$ ตามสมการ~\eqref{eq:sampling_criterion}:
\begin{equation}
    T_s \leq \frac{2\pi}{20\,\omega_n} = \frac{1}{20\,f_n}
    \label{eq:ts_criterion_applied}
\end{equation}

\subsubsection{Discretization ด้วย Euler Forward}

\begin{equation}
    A_d = I + A_c T_s =
    \begin{bmatrix} 1 & T_s \\ -\dfrac{K}{m}T_s & 1 - \dfrac{C}{m}T_s \end{bmatrix}
    \label{eq:Ad_euler}
\end{equation}

\begin{equation}
    B_d = \mathbf{0}_{2\times1} \qquad (u = 0 \text{ สำหรับ Free Vibration})
    \label{eq:Bd_zero}
\end{equation}

เมื่อแทนค่าตัวเลข:
\begin{equation}
    A_d = \begin{bmatrix}
              [\text{ค่าจริง}] & [\text{ค่าจริง}] \\
              [\text{ค่าจริง}] & [\text{ค่าจริง}]
          \end{bmatrix}
    \label{eq:Ad_numerical}
\end{equation}

\subsubsection{สมการ 2-State Discrete Kalman Filter ทั้ง 5 สมการ}

กำหนด $\hat{\mathbf{x}}_k = [\hat{x}_1(k),\; \hat{x}_2(k)]^T$
= $[\text{estimated position},\; \text{estimated velocity}]^T$

\paragraph{ขั้นตอน Prediction (A Priori)}

\begin{equation}
    \hat{\mathbf{x}}_k^{(-)} = A_d\,\hat{\mathbf{x}}_{k-1}
    \tag{2-State-1: State Prediction}
    \label{eq:kf2_predict_state}
\end{equation}

\begin{equation}
    P_k^{(-)} = A_d\,P_{k-1}\,A_d^T + Q
    \tag{2-State-2: Covariance Prediction}
    \label{eq:kf2_predict_cov}
\end{equation}

\paragraph{ขั้นตอน Update (A Posteriori)}

\begin{equation}
    \mathbf{K}_k = P_k^{(-)} H^T \!\left(H P_k^{(-)} H^T + R\right)^{-1}
    \tag{2-State-3: Kalman Gain}
    \label{eq:kf2_gain}
\end{equation}

\begin{equation}
    \hat{\mathbf{x}}_k = \hat{\mathbf{x}}_k^{(-)}
                        + \mathbf{K}_k\!\left(z_k - H\hat{\mathbf{x}}_k^{(-)}\right)
    \tag{2-State-4: State Update}
    \label{eq:kf2_update_state}
\end{equation}

\begin{equation}
    P_k = \left(I - \mathbf{K}_k H\right) P_k^{(-)}
    \tag{2-State-5: Covariance Update}
    \label{eq:kf2_update_cov}
\end{equation}

ขนาดของ Matrix และ Vector ในระบบนี้:

\begin{table}[htbp]
    \centering
    \caption{ขนาดของ Matrix และ Vector ใน 2-State KF}
    \label{tab:matrix_dims}
    \begin{tabular}{clcl}
        \hline
        \textbf{สัญลักษณ์} & \textbf{ชื่อ} & \textbf{ขนาด} & \textbf{หน่วย} \\
        \hline
        $\hat{\mathbf{x}}_k$ & State Estimate     & $2\times1$ & [m, m/s] \\
        $A_d$                & State Transition   & $2\times2$ & --- \\
        $P_k$                & Error Covariance   & $2\times2$ & --- \\
        $Q$                  & Process Noise      & $2\times2$ & --- \\
        $H$                  & Observation Matrix & $1\times2$ & --- \\
        $\mathbf{K}_k$       & Kalman Gain        & $2\times1$ & --- \\
        $R$                  & Measurement Noise  & Scalar     & cm$^2$ \\
        $z_k$                & Measurement        & Scalar     & cm \\
        \hline
    \end{tabular}
\end{table}

%% ─────────────────────────────────────────────────────────────────────────────
\subsection{ขั้นตอนการทดลอง}

\subsubsection{ส่วน A --- Derive และ Discretize Model}

\begin{enumerate}
    \item เขียน Continuous State-Space Model แทนค่า $K$, $C$, $m$ จริง
          คำนวณ $\omega_n$ และ $f_n$ ตามสมการ~\eqref{eq:omega_n}
    \item เลือก $T_s$ ที่ตอบสนองเงื่อนไข $T_s \leq 1/(20 f_n)$
          และสอดคล้องกับ Sampling Time ของ HC-SR04
    \item คำนวณ $A_d$ ด้วย Euler Forward ตามสมการ~\eqref{eq:Ad_euler}
          ตรวจสอบว่า Eigenvalue ของ $A_d$ อยู่ภายใน Unit Circle (Stable)
    \item เขียนสมการ KF ทั้ง 5 สมการในรูปแบบ Matrix ตาม
          สมการ~\eqref{eq:kf2_predict_state}--\eqref{eq:kf2_update_cov}
\end{enumerate}

\subsubsection{ส่วน B --- การกำหนดค่า $R$ และ Matrix $Q$ สำหรับ 2-State KF}

\paragraph{Measurement Noise Covariance $R$}
ค่า $R$ สำหรับ 2-State KF ใช้ค่าเดียวกับที่คำนวณได้จากการทดลองที่~2
ที่ระยะ 20~cm ซึ่งสอดคล้องกับ Operating Range ของระบบ MSD

\begin{equation}
    R = 0.0337\ \text{cm}^2
    \label{eq:R_exp5_measured}
\end{equation}

เหตุผลที่เลือกค่าจากระยะ 20~cm แทนระยะอื่น เนื่องจากมวลในระบบ
MSD เคลื่อนที่อยู่ในช่วงประมาณ $x_\text{eq} \pm 7$~cm
ทำให้ระยะเฉลี่ยระหว่าง Sensor กับมวลอยู่ที่ประมาณ 11--18~cm
ซึ่งใกล้เคียงกับระยะ 20~cm มากที่สุดในชุดการวัดที่มี

\paragraph{Process Noise Matrix $Q$}
กำหนด $Q$ เป็น Diagonal Matrix ขนาด $2\times2$

\begin{equation}
    Q = \begin{bmatrix} q_1 & 0 \\ 0 & q_2 \end{bmatrix}
    \label{eq:Q_matrix}
\end{equation}

การเลือกค่า $q_1$ และ $q_2$ อาศัยหลักการดังนี้

\begin{itemize}
    \item $q_1$ (Position Process Noise) ตั้งค่าน้อยกว่า $q_2$
    เนื่องจาก Position เปลี่ยนแปลงช้า (เป็น Integral ของ Velocity)
    และ Dynamic Model อธิบาย Position ได้แม่นยำในช่วง Displacement เล็ก
    เลือก $q_1 = 1\times10^{-5}$~cm$^2$ ซึ่งมีขนาดใกล้เคียงกับ
    Nominal Value ที่ Al~Tahtawi (2018) ใช้และสอดคล้องกับ
    Variance ของ Sensor ที่ระยะใกล้

    \item $q_2$ (Velocity Process Noise) ตั้งค่าสูงกว่า $q_1$
    เนื่องจาก Velocity ได้รับผลกระทบจากแรงที่ Linear Model
    ไม่ครอบคลุม เช่น Friction Non-linear จาก Guide Rail
    ซึ่งทำให้ Model Prediction ของ Velocity มีความคลาดเคลื่อนสะสมสูงกว่า
    เลือก $q_2 = 1\times10^{-3}$~cm$^2$
    ซึ่งใหญ่กว่า $q_1$ สองลำดับ สะท้อนความไม่แน่นอนที่สูงกว่า
\end{itemize}

ดังนั้นค่าเริ่มต้นของ Matrix $Q$ ที่เลือกใช้คือ

\begin{equation}
    Q = \begin{bmatrix} 1\times10^{-5} & 0 \\ 0 & 1\times10^{-3} \end{bmatrix}
    \text{ cm}^2
    \label{eq:Q_values}
\end{equation}

\subsubsection{ส่วน C --- Implement 2-State KF บน STM32}

\begin{enumerate}
    \setcounter{enumi}{6}
    \item Implement 2-State KF ใน C โดยใช้ Library
          \texttt{kalman.c}/\texttt{kalman.h} ที่พัฒนาไว้
          (ใช้ได้ทันทีโดยตั้งค่า $Q_1 \neq 0$ และ $\mathtt{p11} > 0$
          ต่างจากการทดลองที่ 2 ที่ Collapse เป็น Scalar)
    \item ใน Main Loop: อ่าน Sensor $\rightarrow$ Prediction Step
          $\rightarrow$ Update Step $\rightarrow$
          ส่ง $\hat{x}_1$ (Position) และ $\hat{x}_2$ (Velocity) ออก Serial
\end{enumerate}

\begin{lstlisting}[language=C,
    caption={kalman.h --- 2-State Discrete Kalman Filter — Public Interface},
    label={lst:kalmanh_exp5}]
#ifndef KALMAN_H
#define KALMAN_H

typedef struct {
    float dt;   // Sampling period
    float q0, q1, R; // Noise parameters
    float x0, x1; // State estimates
    float p00, p01, p11; // Error covariance
} Kalman2State_t;

void  Kalman2_Init(Kalman2State_t *kf, float dt, float q0, float q1, float R, float pos, float vel);
float Kalman2_Update(Kalman2State_t *kf, float measurement);
float Kalman2_GetVelocity(const Kalman2State_t *kf);

#endif
\end{lstlisting}

\begin{lstlisting}[language=C,
    caption={kalman.c --- 2-State Discrete Kalman Filter Implementation},
    label={lst:kalmanc_exp5}]
#include "kalman.h"

void Kalman2_Init(Kalman2State_t *kf, float dt, float q0, float q1, float R, float pos, float vel) {
    kf->dt = dt; kf->q0 = q0; kf->q1 = q1; kf->R = R;
    kf->x0 = pos; kf->x1 = vel;
    kf->p00 = 1.0f; kf->p01 = 0.0f; kf->p11 = 1.0f;
}

float Kalman2_Update(Kalman2State_t *kf, float z) {
    // Prediction
    float x0_p = kf->x0 + kf->dt * kf->x1;
    float p00_p = kf->p00 + kf->dt * (2.0f * kf->p01 + kf->dt * kf->p11) + kf->q0;
    float p01_p = kf->p01 + kf->dt * kf->p11;
    float p11_p = kf->p11 + kf->q1;

    // Gain & Update
    float K0 = p00_p / (p00_p + kf->R);
    float K1 = p01_p / (p00_p + kf->R);
    float y = z - x0_p;

    kf->x0 = x0_p + K0 * y;
    kf->x1 = kf->x1 + K1 * y;
    kf->p00 = (1.0f - K0) * p00_p;
    kf->p01 = (1.0f - K0) * p01_p;
    kf->p11 = p11_p - K1 * p01_p;

    return kf->x0;
}

float Kalman2_GetVelocity(const Kalman2State_t *kf) { return kf->x1; }
\end{lstlisting}

\begin{lstlisting}[language=C,
    caption={main.c --- 2-State KF สำหรับระบบ MSD (Free Vibration)},
    label={lst:main_exp5}]
// System & Noise Parameters
#define KALMAN_DT     0.010f
#define KALMAN_Q0     [q1]
#define KALMAN_Q1     [q2]
#define KALMAN_R      0.0337f

HAL_TIM_Base_Start(&htim1);
HCSR04_Init(&htim1);

float raw_cm = HCSR04_Read();
Kalman2_Init(&kf_msd, KALMAN_DT, KALMAN_Q0, KALMAN_Q1, KALMAN_R, raw_cm, 0.0f);

while (1) {
    raw_cm = HCSR04_Read();
    if (raw_cm > 0.0f) {
        float pos = Kalman2_Update(&kf_msd, raw_cm);
        float vel = Kalman2_GetVelocity(&kf_msd);
        printf("%lu,%.4f,%.4f,%.4f\r\n", HAL_GetTick(), raw_cm, pos, vel);
        fflush(stdout);
    }
    HAL_Delay(10);
}
\end{lstlisting}

\subsubsection{ส่วน D --- ทดสอบ Free Vibration}

\begin{enumerate}
    \setcounter{enumi}{8}
    \item กำหนด $x = 0$ ที่ Equilibrium ดึงมวลไปที่ $x_0 = 2$~cm แล้วปล่อย
          บันทึก 15 วินาที (Raw $z_k$, Position Estimate, Velocity Estimate)
    \item ทำซ้ำที่ $x_0 = 4$~cm และ $x_0 = 6$~cm
    \item Simulate MSD Model จากการทดลองที่ 4 ใน MATLAB
          ด้วย Initial Condition เดียวกัน เปรียบเทียบกับ KF Output
\end{enumerate}

\subsubsection{ส่วน E --- Sensitivity Study}

\begin{enumerate}
    \setcounter{enumi}{11}
    \item ที่ $x_0 = 4$~cm ทดสอบ $Q$ อย่างน้อย 3 ชุด
          (Scale Diagonal ด้วย $0.01\times$, $1\times$, $100\times$)
    \item ปรับ $R$ ($0.5\times$, $1\times$, $2\times$ ของ Measured Variance)
          บันทึกการเปลี่ยนแปลงของ Position Estimate
\end{enumerate}

%% ─────────────────────────────────────────────────────────────────────────────
\subsection{ผลการทดลอง}

\subsubsection{Continuous State-Space Matrices}

\begin{equation}
    A_c = \begin{bmatrix} 0 & 1 \\ -K/m & -C/m \end{bmatrix}
        = \begin{bmatrix} 0 & 1 \\ [\text{ค่าจริง}] & [\text{ค่าจริง}] \end{bmatrix},
    \quad
    B_c = \begin{bmatrix} 0 \\ 1/m \end{bmatrix}
        = \begin{bmatrix} 0 \\ [\text{ค่าจริง}] \end{bmatrix},
    \quad
    H = \begin{bmatrix} 1 & 0 \end{bmatrix}
    \label{eq:cont_matrices_numerical}
\end{equation}

\subsubsection{Sampling Rate และ Discrete State-Space Matrices}

\begin{equation}
    \omega_n = \sqrt{\frac{K}{m}} = [\text{ค่าจริง}] \; \text{rad/s},
    \qquad
    f_n = \frac{\omega_n}{2\pi} = [\text{ค่าจริง}] \; \text{Hz}
    \label{eq:omega_n_numerical}
\end{equation}

\begin{equation}
    T_s \leq \frac{1}{20 f_n} = [\text{ค่าจริง}] \; \text{ms}
    \quad \Rightarrow \quad \text{เลือก } T_s = 10 \; \text{ms}
    \label{eq:ts_choice}
\end{equation}

\begin{equation}
    A_d = I + A_c T_s
        = \begin{bmatrix} 1 & T_s \\ -\frac{K}{m}T_s & 1-\frac{C}{m}T_s \end{bmatrix}
        = \begin{bmatrix} [\text{ค่าจริง}] & [\text{ค่าจริง}] \\
                          [\text{ค่าจริง}] & [\text{ค่าจริง}] \end{bmatrix}
    \label{eq:Ad_result}
\end{equation}

\begin{equation}
    B_d = \mathbf{0}_{2\times1}
    \qquad (u = 0 \text{ สำหรับ Free Vibration})
    \label{eq:Bd_result}
\end{equation}

\begin{table}[htbp]
    \centering
    \caption{ค่าพารามิเตอร์ KF ที่ใช้ในการทดสอบ}
    \label{tab:kf2_params}
    \begin{tabular}{lccc}
        \hline
        \textbf{พารามิเตอร์} & \textbf{สัญลักษณ์} & \textbf{หน่วย} & \textbf{ค่า} \\
        \hline
        Position Process Noise  & $q_1$  & ---         & [ค่าจริง] \\
        Velocity Process Noise  & $q_2$  & ---         & [ค่าจริง] \\
        Measurement Noise       & $R$    & cm$^2$      & $0.0337$  \\
        Sampling Period         & $T_s$  & ms          & $10$      \\
        Initial Position Estimate & $\hat{x}_1(0)$ & cm & $z_0$ (การวัดแรก) \\
        Initial Velocity Estimate & $\hat{x}_2(0)$ & cm/s & $0$ \\
        Initial Covariance      & $P_0$  & ---         & $I_{2\times2}$ \\
        \hline
    \end{tabular}
\end{table}

\begin{figure}[htbp]
    \centering
    % \includegraphics[width=0.85\textwidth]{images/exp5_x0_2cm.png}
    \vspace{5cm}
    \caption{Free Vibration ที่ $x_0 = 2$~cm: Raw $z_k$ (เส้นเทา),
             KF Position Estimate (เส้นน้ำเงิน), และ Mathematical Model Simulation
             (เส้นประแดง) บนแกนเดียวกัน (แกน $x$: เวลา (s), แกน $y$: ระยะทาง (m))}
    \label{fig:exp5_2cm}
\end{figure}

\begin{figure}[htbp]
    \centering
    % \includegraphics[width=0.85\textwidth]{images/exp5_x0_4cm.png}
    \vspace{5cm}
    \caption{Free Vibration ที่ $x_0 = 4$~cm: Raw $z_k$, KF Position Estimate,
             และ Mathematical Model Simulation}
    \label{fig:exp5_4cm}
\end{figure}

\begin{figure}[htbp]
    \centering
    % \includegraphics[width=0.85\textwidth]{images/exp5_x0_6cm.png}
    \vspace{5cm}
    \caption{Free Vibration ที่ $x_0 = 6$~cm: Raw $z_k$, KF Position Estimate,
             และ Mathematical Model Simulation}
    \label{fig:exp5_6cm}
\end{figure}

\begin{figure}[htbp]
    \centering
    % \includegraphics[width=0.85\textwidth]{images/exp5_velocity_estimate.png}
    \vspace{5cm}
    \caption{KF Velocity Estimate $\hat{x}_2$ (เส้นน้ำเงิน) เทียบกับ
             Finite Difference ของ Position
             $v_\text{fd}(k) = (z_k - z_{k-1})/T_s$ (เส้นเทา)
             แสดงให้เห็นว่า KF Velocity เรียบกว่าแต่มีแนวโน้มเดียวกัน}
    \label{fig:exp5_velocity}
\end{figure}

\begin{figure}[htbp]
    \centering
    % \includegraphics[width=0.85\textwidth]{images/exp5_Q_sensitivity.png}
    \vspace{5cm}
    \caption{Sensitivity Study: KF Position Output สำหรับ $Q$ ที่ Scale ด้วย
             $0.01\times$, $1\times$, $100\times$ บนแกนเดียวกัน ที่ $x_0 = 4$~cm}
    \label{fig:exp5_Q_sensitivity}
\end{figure}

\begin{table}[htbp]
    \centering
    \caption{ผลการเปรียบเทียบ KF กับ Mathematical Model สำหรับระยะ Release ต่างๆ}
    \label{tab:exp5_comparison}
    \begin{tabular}{cccc}
        \hline
        \textbf{ระยะ $x_0$}
            & \textbf{Matrix $Q$}
            & \textbf{ค่า $R$ (cm$^2$)}
            & \textbf{การเปรียบเทียบ KF vs.\ Model} \\
        \hline
        2~cm & $\text{diag}(q_1, q_2)$ & $0.0337$ & [เชิงคุณภาพ] \\
        4~cm & $\text{diag}(q_1, q_2)$ & $0.0337$ & [เชิงคุณภาพ] \\
        6~cm & $\text{diag}(q_1, q_2)$ & $0.0337$ & [เชิงคุณภาพ] \\
        \hline
    \end{tabular}
\end{table}

%% ─────────────────────────────────────────────────────────────────────────────
\subsection{สรุปผลการทดลอง}

\begin{itemize}
    \item 2-State KF สามารถประมาณค่าทั้ง Position และ Velocity
          โดย Position Estimate เรียบกว่า Raw Data และ Velocity Estimate
          สอดคล้องทางกายภาพกับ Derivative ของ Displacement
    \item KF และ Mathematical Model ตรงกันได้ดีที่ Displacement เล็ก
          ($x_0 = 2$~cm) แต่ที่ Displacement ใหญ่ ($x_0 = 6$~cm)
          อาจมีความเบี่ยงเบนเนื่องจาก Linear Model
          มีความแม่นยำลดลงนอกช่วง Small Oscillation
    \item การเพิ่ม $Q$ ทำให้ KF ติดตาม Measurement มากขึ้น
          (สัญญาณออก Noisy กว่า) ส่วนการเพิ่ม $R$ ทำให้ KF
          พึ่ง Model Prediction มากขึ้น (สัญญาณออกเรียบกว่าแต่ช้ากว่า)
    \item Velocity Estimate สมเหตุสมผลทางกายภาพ:
          $\hat{x}_2 \approx 0$ ที่จุด Maximum Displacement
          และ $|\hat{x}_2|$ สูงสุดที่ Equilibrium Crossing
          สอดคล้องกับ Energy Conservation
    \item ข้อจำกัดสำคัญ: Noise Floor ของ HC-SR04 ($\approx 2$~cm),
          ข้อสมมติ Linear Model, และ Non-Adaptive Filter
          ที่ Q กับ R คงที่ตลอดเวลา
\end{itemize}

%% ─────────────────────────────────────────────────────────────────────────────
\subsection{อภิปรายผล}

\subsubsection{ความหมายทางกายภาพของ Matrix $Q$ และ Scalar $R$}

\begin{itemize}
    \item \textbf{$Q_{11}$ (Position Process Noise):}
          แสดงความไม่แน่นอนใน Prediction ของ Position
          ค่านี้มักเล็กเนื่องจาก Position เปลี่ยนแปลงช้า
          (ต้อง Integrate จาก Velocity) และ Dynamic Model
          อธิบาย Position ได้แม่นยำในช่วง Displacement เล็ก

    \item \textbf{$Q_{22}$ (Velocity Process Noise):}
          แสดงความไม่แน่นอนใน Prediction ของ Velocity
          ค่านี้ใหญ่กว่า $Q_{11}$ เพราะ Velocity ได้รับผลกระทบ
          จากแรงที่ Linear Model ไม่ครอบคลุม เช่น แรงเสียดทาน
          Non-linear และ Air Resistance

    \item \textbf{Off-Diagonal Elements = 0:}
          สมมติว่า Process Noise ของ Position และ Velocity
          ไม่มี Correlation กัน ซึ่งลดความซับซ้อนในการ Tune
          และเป็น Conservative Assumption ที่ยอมรับได้
    
    \item คำถามที่น่าพิจารณาคือเหตุใด Off-diagonal ของ $Q$ จึงตั้งเป็นศูนย์
            แม้ว่า Friction จาก Guide Rail จะกระทำต่อ Velocity State โดยตรง
            เหตุผลคือ Friction เป็น Deterministic Force ที่ขึ้นกับความเร็ว
            ไม่ใช่ Stochastic Process Noise ที่เป็น Random ดังนั้น
            ผลของ Friction จึงถูก Absorb เข้าไปใน $q_2$ (Velocity Process Noise)
            แทนที่จะสร้าง Correlation ระหว่าง Position และ Velocity Noise
            การตั้ง Off-diagonal เป็นศูนย์จึงเป็น Conservative Assumption
            ที่ยอมรับได้และลดความซับซ้อนในการ Tune พารามิเตอร์
            หาก Friction มีลักษณะ Stochastic อย่างแท้จริง
            การใช้ Cross-Covariance ที่ไม่เป็นศูนย์อาจให้ผลดีกว่า
            แต่ต้องการข้อมูลเพิ่มเติมในการประมาณค่า

    \item \textbf{$R$ (Measurement Noise):}
          ได้จากการวัด Variance ของ HC-SR04 จริงในการทดลองที่ 2
          ค่า $R = 0.0337$~cm$^2$ บอกว่า Filter ควรเชื่อ Measurement
          น้อยลงเมื่อ Sensor มี Noise สูง ซึ่ง Balance กับ $Q$
          ในการกำหนดสัดส่วนความเชื่อมั่น
\end{itemize}

\subsubsection{เหตุผลที่ใช้ Displacement จาก Equilibrium เป็น State}

เมื่อใช้ Equilibrium เป็นจุดอ้างอิง แรงโน้มถ่วง $mg$ ถูก Balance
โดย Spring Pre-load $Kx_\text{eq} = mg$ พอดี ทำให้ EOM ไม่มี Constant Term:
\begin{equation}
    m\ddot{x} + C\dot{x} + Kx = 0
\end{equation}
ส่งผลให้ State-Space Model มีรูปแบบ $\dot{\mathbf{x}} = A_c\mathbf{x}$
ที่สะอาดโดยไม่มี Bias Term ใน Prediction Step
และ $u = 0$ ตลอดการทดลอง Free Vibration

\subsubsection{การวิเคราะห์ Velocity Estimate}

ตรวจสอบความสมเหตุสมผลของ Velocity Estimate ด้วย Energy Conservation:
จาก $\frac{1}{2}Kx_0^2 = \frac{1}{2}mv_\text{max}^2$ จะได้:

\begin{equation}
    v_\text{max} = x_0\sqrt{\frac{K}{m}} = x_0\,\omega_n
    \label{eq:v_max_theory}
\end{equation}

KF Velocity Estimate สมเหตุสมผลถ้า:
\begin{itemize}
    \item $|\hat{x}_2|_\text{max} \approx v_\text{max}$ ตามสมการ~\eqref{eq:v_max_theory}
    \item $\hat{x}_2 \approx 0$ ที่จุด Maximum Displacement (Peak ของ $z_k$)
    \item $\hat{x}_2$ มีเครื่องหมายสลับกันทุกครึ่ง Period
          สอดคล้องกับการแกว่งของระบบ
    \item แนวโน้มของ $\hat{x}_2$ ตรงกับ Finite Difference
          $v_\text{fd}(k) = (z_k - z_{k-1})/T_s$ แม้ Finite Difference
          จะมี Noise สูงกว่า

\end{itemize}

\subsubsection{เปรียบเทียบ KF กับ Mathematical Model}

\begin{itemize}
    \item \textbf{ที่ $x_0 = 2$~cm (Small Oscillation):}
          Linear Model ควรแม่นยำที่สุด เนื่องจากสปริงและ Damper
          ทำงานในช่วง Linear และ EOM เป็นสมการที่ถูกต้อง
          KF Position Estimate ควรใกล้เคียง Model Simulation มาก

    \item \textbf{ที่ $x_0 = 6$~cm (Larger Oscillation):}
          อาจพบความเบี่ยงเบนเนื่องจาก:
          \begin{itemize}
              \item สปริงทำงานนอกช่วง Linear Elastic บางส่วน
              \item Damping ที่แท้จริงอาจไม่คงที่ (Velocity-dependent Friction)
              \item HC-SR04 Noise Profile อาจเปลี่ยนแปลงตามระยะ
          \end{itemize}
          KF ยังคง Estimate ได้สมเหตุสมผลเพราะมีค่า $Q$ ที่รองรับ
          ความไม่แน่นอนของ Model แต่ความแตกต่างจาก Simulation
          จะมากกว่ากรณี $x_0 = 2$~cm
\end{itemize}

\subsubsection{ข้อจำกัดของระบบและแนวทางพัฒนา}

\begin{itemize}
    \item \textbf{Noise Floor ของ HC-SR04:} เซนเซอร์วัดได้ต่ำสุดประมาณ 2~cm
          ดังนั้น Oscillation ที่ Amplitude เล็กมากอาจถูก Noise กลบ
          ทำให้ KF ไม่สามารถประมาณ State ได้แม่นยำในช่วงท้าย
          ของ Free Vibration ที่ Amplitude ลดลงมาก

    \item \textbf{ข้อสมมติ Linear Model:} EOM สมมติว่า $K$ และ $C$
          คงที่และเป็น Linear ซึ่งอาจไม่สมบูรณ์ที่ Amplitude ใหญ่
          หากต้องการความแม่นยำสูงขึ้น ควรใช้ Extended Kalman Filter (EKF)
          กับ Nonlinear Model

    \item \textbf{Non-Adaptive Filter:} KF นี้ใช้ $Q$ และ $R$ คงที่ตลอดเวลา
          ซึ่งอาจไม่เหมาะสมหาก Noise Characteristics เปลี่ยนแปลง
          การพัฒนาต่อเป็น Adaptive KF จะช่วยให้ระบบ Tune ตัวเองได้

    \item \textbf{ความคลาดเคลื่อนจาก Discretization:}
          Euler Forward มี Truncation Error $O(T_s^2)$
          ที่ $T_s = 10$~ms ความต่างจาก ZOH ประมาณ 2\%
          ซึ่งยอมรับได้สำหรับ Application นี้แต่ควรระบุไว้ใน Limitation
\end{itemize}